\begin{question}
\noindent $O$ is the centre of the circle below. $A$ and $B$ are points on the circumference, and $OA$ and $OB$ are radii of the circle. $A$ and $B$ are joined to form the triangle $OAB$.

\begin{center}
\begin{tikzpicture}[scale=2.4]
    \coordinate (O) at (0,0);
    \coordinate (A) at (110:1);
    \coordinate (B) at (70:1);

    \draw[thick] (0,0) circle (1);
    \fill[black] (O) circle (0.4pt);

    \draw[thick] (O) -- (A) node[midway, left] {};
    \draw[thick] (O) -- (B) node[midway, right] {};
    \draw[thick] (A) -- (B);

    \node[below] at (O) {$O$};
    \node[above left] at (A) {$A$};
    \node[above right] at (B) {$B$};

    \pic [draw, angle radius=6mm, "$40^\circ$"] {angle = B--O--A};
\end{tikzpicture}
\end{center}

\noindent Angle $AOB = 40^\circ$.

Work out the size of angle $OAB$. Give a reason for your answer.

\vspace{4cm}

\begin{flushright}
    \answerequnit{OAB}{$^\circ$}{2}
\end{flushright}
\end{question}
